\section{Вероятностная модель ARC-AGI}

\subsection{Решетки и скрытое правило}

Побочная задача ARC-AGI используется не как основной источник результатов, а как проверка переносимости идеи времени остановки. В этой задаче вход представляет собой цветную решетку
\begin{equation}
  A\in\{1,\ldots,C\}^{H\times W},
\end{equation}
а ответ --- выходную решетку
\begin{equation}
  Y\in\{1,\ldots,C\}^{H'\times W'}.
\end{equation}
Для одной задачи предполагается скрытое правило преобразования \(\Theta\), по которому несколько обучающих пар согласованы между собой. Условно это можно записать как
\begin{equation}
  Y=f_{\Theta}(A).
  \label{eq:arc_rule}
\end{equation}

Рекурсивная модель в этом случае строит последовательность распределений по цветам клеток:
\begin{equation}
  Q_k(i,j,c)=\Prob(Y_{i,j}=c\mid A,\Hh_k), \qquad c=1,\ldots,C.
  \label{eq:arc_cell_distribution}
\end{equation}
Потеря всей решетки при остановке на шаге \(k\) задается суммой по клеткам:
\begin{equation}
  L_k^{\mathrm{arc}}=-\sum_{i,j}\log Q_k(i,j,Y_{i,j}).
  \label{eq:arc_loss}
\end{equation}
Оракульное время остановки определяется так же, как в языковой задаче:
\begin{equation}
  \tau_{\mathrm{arc}}^\star=\argmin_{0\le j\le K}\{L_j^{\mathrm{arc}}+\lambda c(j)\}.
  \label{eq:arc_tau}
\end{equation}

Для ARC-AGI естественно использовать одно время остановки на всю решетку. Отдельное время для каждой клетки технически возможно, но нарушает смысл задачи: ответ оценивается как единый объект, а не как независимый набор клеток. Поэтому выход распределительной модели имеет размер
\begin{equation}
  B\times (K+1)\times (K+1),
  \label{eq:arc_output_size}
\end{equation}
где \(B\) --- число задач или примеров в пакете.

\subsection{Роль побочной задачи}

В языковой задаче трудность позиции может быть связана с грамматикой, частотностью словосочетания, дальними зависимостями и неоднозначностью смысла. В ARC-AGI трудность имеет другую природу: модель должна восстановить скрытое преобразование решетки. Если распределение времени остановки в двух задачах демонстрирует сходные режимы, это поддерживает гипотезу, что предсказание времени остановки описывает общую вычислительную сложность объекта, а не только частный эффект языковых данных.

В данной работе ARC-AGI не должен раздувать постановку. Он используется в трех местах: как побочная модель, как отдельная строка в экспериментальном плане и как проверка правил остановки на решеточных данных.

